\documentclass[10pt]{report}
%\usepackage{}
\usepackage[T1]{fontenc}
\usepackage[latin1]{inputenc}
%\usepackage[francais]{babel}
\usepackage{a4}
\usepackage{fullpage}
\usepackage{fancyheadings}
\pagestyle{empty}
\usepackage{epsfig}
\setlength{\headrulewidth}{0pt}
\setlength{\footskip}{0.5cm}
\usepackage{anysize}
\marginsize{2.5cm}{2cm}{2cm}{2cm}
\usepackage{url} \urlstyle{tt}


\begin{document}


\thispagestyle{fancy}

\chead{
\ \\[0.5cm]
\makebox[0pt]{
\begin{tabular}{rcr}
\begin{minipage}{25mm}
\centerline{\psfig{file=LORIA.eps,width=25mm}}
\end{minipage}&
\begin{minipage}{10cm}
\begin{center}
{\large\bf Laboratoire Lorrain de Recherche\\ en  Informatique et ses Applications}\\
\noindent UMR N$^\circ$ 7503, CNRS, INRIA\\ UHP-Nancy 1, INPL , Nancy 2
\end{center}
\end{minipage} &
\begin{minipage}{30mm}
\centerline{\psfig{file=INRIA.eps,height=10mm}}
\end{minipage}%\\\hline
\end{tabular}
}}


\cfoot{\small
%hrule
%vspace{2mm}
\makebox[0pt]{
\begin{tabular}{rrrr}
\begin{minipage}{25mm}
\centerline{\psfig{file=CNRS.eps,width=25mm}}
\end{minipage}&
\begin{minipage}{6.6cm}
LORIA, Universit� Henri Poincar�\\
Campus Scientifique - BP 239\\
 54506 Vand{\oe}uvre-l�s-Nancy Cedex, France\end{minipage} &
\begin{minipage}{3.7cm}
Fax :  +33 3 83 59 30 00 \\
T\'el : +33 3 83 27 83 19\\
\url{http://www.loria.fr}
\end{minipage} &
\begin{minipage}{25mm}
\centerline{\psfig{file=logo_nu,width=40mm}}
\end{minipage}
\end{tabular}
}
\ \\[5cm]
}

\newcommand{\edff}[1]{\mbox{\sf EDF$^{(#1)}$}}
\newcommand{\edf}{\mbox{\sf EDF}}

\ \\[2cm]

\hfill
\begin{tabular}{l}
\smallskip\\
Emmanuel Jeannot\\
Charg� de recherche INRIA\\
Tel : +33 3 83 54 95 85 11\\
{\tt emmanuel.jeannot@loria.fr}\\[5mm]

Nancy,   \today
\end{tabular}
\ \\[1cm]
\noindent{� l'agent comptable de l'UHP}

J'ai eu la chance d'encadrer le stage de licence d'Eloi du Bois du 25 juin au 25 ao�t 2007 intitul� "{\it compression avec perte d'images � la vol�e dans AdOC}".  Il m'avait �t� chaudement recommand� par C�cile Germain, Professeur des Universit�s � Orsay qui avait d�j� travaill� avec lui sur la compression d'images par ondelettes, th�me en relation directe avec le sujet de stage. Il avait donc une  comp�tence tr�s rare � ce niveau de formation. Quand je lui ai propos� de venir en stage � Nancy, la question financi�re c'est tout de suite pos�e. En effet, il s'agit d'un �tudiant ayant de faibles ressources et qui n'a pas de famille sur Nancy. Ainsi, comme il m'�tait recommand� et qu'il avait les comp�tences requises j'ai d�cider de lui payer son logement pour les deux mois de stage sur les cr�dits de l'ACI masses de donn�es AGIR (dont le sujet du stage est un des th�mes). 

Au final, les bonnes conditions de ce stage lui on permit d'effectuer un excellent travail autant d'un point de vue qualitatif que quantitatif.   Ses contributions sont d'un tr�s bon niveau et feront prochainement l'objet d'une publication de niveau international. En outre, les d�veloppements effectu�s,  seront int�gr�s dans notre biblioth�que de compression adaptative AdOC.

En conclusion, Eloi Du Bois s'est av�r� un tr�s bon stagiaire. Si je n'avais pas pu lui payer son logement sur Nancy, il n'aurait pas pu venir et mes travaux sur la compression adaptative dans l'ACI AGIR auraient �t�, au mieux, retard�s, au pire inachev�s. D'un point de vue scientifique, sa venue a donc �t� tr�s bon investissement. 

Bien cordialement, 
\hfill{Emmanuel Jeannot}\\[2mm]

%\hspace{12cm}\includegraphics[width=4cm]{signature.ps}

\vspace{.2in}

\end{document}



